\chapter{Introduction}
\thispagestyle{plain}

The classification of cloud computing requirements from natural language documents is a well-studied problem in requirements engineering. Casalicchio and Cotumaccio's AI-CRAS framework \cite{casalicchio2024aicras} addressed it by fine-tuning a BERT model on a domain-specific dataset of cloud requirement sentences, achieving a recall of 0.96 and an F1-score of 0.92 on the binary task of identifying whether a sentence constitutes a requirement.

What AI-CRAS does not examine is the effect of \textit{ambiguity}. Many requirements are phrased in ways that allow more than one interpretation — due to vague wording, grammatical structure, or missing context. Natural language requirements can exhibit several distinct forms of ambiguity: lexical, syntactic, semantic, pragmatic, and language-error (described in detail in Chapter~3, following the taxonomy of \cite{bano2015ambiguity}). These are common in real-world requirement documents, but their effect on model performance has not been studied.

This thesis tests whether ambiguity matters for classification — and whether certain categories are more vulnerable to it than others. Using the same AI-CRAS dataset, each sentence is manually annotated with ambiguity labels. To answer RQ1, a BERT classifier is trained on the full dataset (baseline) and on the unambiguous subset and the two are compared. To answer RQ2, the same classifier is re-run on an LLM-deambiguified version of the dataset. To answer RQ3, an LLM is prompted to classify the sentences directly, in both zero-shot and few-shot configurations. For each experiment, results are broken down by the six cloud service categories to assess whether ambiguity affects them uniformly.

The three research questions are:

\begin{itemize}
    \item \textbf{RQ1:} Does ambiguity negatively affect BERT classification performance, and which of the six cloud service categories is most affected?
    \item \textbf{RQ2:} Can LLM-based rewriting of ambiguous sentences improve BERT's accuracy — both overall and for the most ambiguity-sensitive categories?
    \item \textbf{RQ3:} Can an LLM used directly as a classifier match or exceed BERT?
\end{itemize}

The main contributions of this thesis are: an ambiguity-annotated extension of the AI-CRAS dataset; an empirical analysis of how ambiguity affects BERT-based classification, including per-category breakdowns; and a comparative evaluation of LLMs as both a disambiguation tool and a direct classifier.
